\documentclass[manuscript, screen]{jair}

\setcopyright{none}
\acmDOI{}
\JAIRAE{}
\JAIRTrack{}
\acmVolume{}
\acmArticle{}
\acmMonth{10}
\acmYear{2026}

\RequirePackage[
  datamodel=acmdatamodel,
  style=acmauthoryear,
  backend=biber,
  giveninits=true,
  uniquename=init
  ]{biblatex}
\renewcommand*{\bibopenbracket}{(}
\renewcommand*{\bibclosebracket}{)}
\addbibresource{references.bib}

\usepackage{tabularx}

\newcommand{\lean}[1]{\texttt{\detokenize{#1}}}
\newcommand{\cls}[1]{\textsc{#1}}
\newcommand{\Sol}{\mathit{Sol}}
\newcommand{\precm}{\mathit{prec}}
\newcommand{\tnI}{\mathit{tn}_I}
\newcolumntype{L}{>{\raggedright\arraybackslash}X}

\begin{document}

\title[Formal Proofs vs.\ the Article]{Where the Lean Formalisation Differs from
  ``Propositional HTN Planning: Problem Classes and Complexity Results''}

\author{Gregor Behnke}
\email{g.behnke@uva.nl}
\affiliation{%
  \institution{University of Amsterdam}
  \city{Amsterdam}
  \country{The Netherlands}}

\renewcommand{\shortauthors}{Behnke}

\begin{abstract}
  All results of the article by \textcite{bercher:et:al:htn-complexity} have been formalised in
  Lean~4 with Mathlib \parencite{moura:ullrich:cade-21,mathlib:cpp-20}. Every theorem depends
  only on the standard axioms \lean{propext}, \lean{Classical.choice} and \lean{Quot.sound}.
  This report lists every place where the formal development deviates from the article. Three
  lemmas state bounds that are false (Lemmas~4.2.4, 4.2.15 and 4.2.21), one lemma misses a case
  (Lemma~4.2.2), one corollary needs $h \ge 1$ (Corollary~4.3.3), and two printed constructions are
  wrong (Theorems~4.3.7 and~4.3.8). Five proofs follow a different route than the article,
  either because the article's step is wrong or imprecise, or because it rests on a cited result
  that we replaced by a self-contained argument. The proofs of all remaining results follow the
  article.
\end{abstract}

\maketitle

\section{Introduction}

The formalisation is spread over four Lean repositories: \lean{complexity-theory} (machines,
languages and complexity classes), \lean{htn-planning} (the HTN model and the problem classes),
\lean{htn-complexity} (the complexity results of the article) and \lean{htn-compilability}. The
article's results live in the namespace \lean{HTN.Article}; the file \lean{ARTICLE.md} in
\lean{htn-complexity} maps every numbered result to its Lean name.

We followed three rules. First, the article's own model is formalised literally and related to
an older, simpler model by bridge lemmas. Second, every complexity result is stated about a
literal language of bit-encoded problems. Third, the article's proof route is always tried
first. When it fails, the error is recorded and repaired, and only when no repair is possible
is a different route used.

Section~\ref{sec:setting} describes how the formal setting differs from the article's.
Section~\ref{sec:statements} lists the results whose statement had to change.
Section~\ref{sec:defs} lists corrections to definitions and constructions that leave the
statements intact. Section~\ref{sec:routes} describes the proofs that take a different route.
Section~\ref{sec:status} gives the complete status.

\section{The Formal Setting}\label{sec:setting}

\paragraph{Model.} The article's model, with negative preconditions and goals, the requirement
that all linearisations of a solution be executable, ordering insertion and task insertion
(TIHTN), is formalised in \lean{HtnPlanning/Article/}. The development predating this work uses
a simpler model with positive preconditions only and a single executable linearisation. Bridge
lemmas connect the two: negative preconditions are compiled away by dual-rail facts, and
ordering insertion corresponds to choosing one executable linearisation. The existing proofs
remain on the old model; the article versions of the theorems are derived through the bridges.

\paragraph{Decomposition trees.} Definitions~2.8--2.10 use an explicit mapping $\beta$ and
isomorphisms of task networks. The formal trees are indexed by the task ids of the applied
method, so $\beta$ and the isomorphisms are implicit. Yields are action sequences.

\paragraph{Missing references.} The draft leaves the references for TIHTN solutions, for the
lemma that acyclic decompositions suffice under task insertion, and for Proposition~2.1.1 as
``??''. We defined TIHTN solutions following \textcite{geier:bercher:ijcai-11} and
\textcite{alford:et:al:icaps-15}, and proved the missing lemma
(\lean{solvableTI_iff_acyc_trees}).

\paragraph{Inputs.} Definition~4.1 speaks of plan-existence problems without fixing an
encoding. We state every result about a language of bit strings that encode a problem
description (domain and instance). This requires recognising the codes of descriptions, and
membership in each problem class, within the resources of the class at hand: polynomial time
for most results, and logarithmic space for Theorem~4.2.10.

\paragraph{Complexity classes.} \cls{P}, \cls{NP} and \cls{PSPACE} are defined via
Cobham-style function classes, which are proved equal to the classes defined by Turing
machines. \cls{EXPTIME} is defined as \cls{P} on exponentially padded inputs. \cls{L} is
defined by deterministic register programs: while programs over a fixed number of registers,
each bounded by a polynomial in the input length and thus holding $O(\log n)$ bits, that read
the input bit by bit. Both are standard characterisations.

\paragraph{Cited results.} The article cites the undecidability of the intersection of
context-free languages, the \cls{PSPACE}-hardness of STRIPS \parencite{bylander:aij-94}, the
\cls{NP}-hardness of SAT \parencite{cook:stoc-71} and the \cls{EXPTIME}-hardness of the game
$G_6$ \parencite{stockmeyer:chandra:siamjc-79}. The first three are proved in the
formalisation: CFG intersection from two-stack machines \parencite{minsky:67}, STRIPS by
Bylander's encoding of space-bounded machines, and SAT by the Cook--Levin construction. The
fourth is avoided; see Section~\ref{sec:routes}.

\section{Results Whose Statement Changed}\label{sec:statements}

\paragraph{Lemma 4.2.4 (number of decompositions).} The article bounds the number of
decompositions of an acyclic problem by $k \cdot m^{h-1} \le r^h$. This bound is false. A
decomposition tree of height $h$ and branching $m$ has up to $1 + m + \dots + m^{h-1}$ inner
nodes, not $m^{h-1}$. A counterexample has $\tnI = \{A, A\}$ and methods $A \to \{B, B\}$ and
$B \to \{a\}$. Here $h = k = m = r = 2$, but a solution needs $2 + 4 = 6 > 4$ decompositions. The
formal lemma \lean{lem_4_2_4} proves the corrected bound
\[
  k \cdot (1 + m + \dots + m^{h-1}) < (r+1)^h .
\]
The uses of the lemma in Theorems~4.2.5 and~4.2.6 and Corollary~4.2.8 only need an exponential
bound and are unaffected.

\paragraph{Lemma 4.2.15 (progression bound for tail-recursive problems).} The bound
$k \cdot m^h$ is $0$ when $m = 0$ and $h \ge 1$, which is wrong. The formal lemma assumes
$m \ge 1$. This is harmless, since any upper bound on method sizes can be taken to be at least~1.

\paragraph{Lemma 4.2.21.} The bound $k \cdot m \cdot h$ fails for $h = 0$ (problems without
compound tasks, where the bound would be $0$) and for $m = 0$. The formal version
\lean{lem_4_2_21_article} proves $k \cdot (1 + (m-1) \cdot h)$. This is at most $k \cdot m
\cdot h$ whenever $m, h \ge 1$, so the article's bound holds in that range.

\paragraph{Lemma 4.2.2 (mostly acyclic to acyclic).} The construction collapses only strata
with at least two elements. A singleton stratum $\{c\}$ with a method $(c, [c])$, which
Definition~3.7 allows, therefore stays cyclic. The formal construction drops every method that
becomes a self-loop (\lean{HTN.Method.IsLoop}).

\paragraph{Proposition 2.1.1 and Corollary 4.2.3.} Both claim an equality of solution sets,
$\Sol(P) = \Sol(P')$. Since decomposition creates fresh task ids, the two sets agree only up to
renaming of ids. The formal statements prove equality up to isomorphism, for HTN and TIHTN
solutions, and equality of plan existence exactly.

\paragraph{Corollary 4.3.3.} The corollary claims \cls{PSPACE}-hardness of
$\cls{TAIL-REC}_h$ for a fixed constant $h$. For $h = 0$ the class has no compound tasks, so it
lies in \cls{NP} (Corollary~4.2.9). The formal statement requires $h \ge 1$.

\section{Corrected Definitions and Constructions}\label{sec:defs}

The following corrections do not change the statement of any result.
Table~\ref{tab:defs} summarises them.

\begin{table}
  \caption{Corrections that leave the statements of the results unchanged.}
  \label{tab:defs}
  \small
  \begin{tabularx}{\linewidth}{@{}lLL@{}}
    \toprule
    Place & In the article & In the formalisation \\
    \midrule
    \S2, $\gamma$ & $\gamma(s, a)$ defined iff $s \subseteq \precm^+(a)$ & $\precm^+(a) \subseteq s$ \\
    Def.~2.5 & Inserting $(t, t')$ with $t' \prec^* t$ or $t = t'$ does not give a strict partial
      order, contrary to the remark after the definition
      & Insertion requires the result to be a strict partial order \\
    Def.~4.2 & Inserted tasks get the ancestors $a(t) \cup \{\alpha(t)\}$, so a name may repeat once
      ($c \to d \to d$) & Formalised literally; harmless \\
    Def.~4.3 & Progression bound quantifies over runs from the initial state only
      & All reachable states \\
    Thm.~4.2.12 & Second item reads $\Sol(P)$ & TIHTN solutions \\
    Thm.~4.3.4 & Bound $2^{\#V}$, $\#V$ not defined & $\#V$ = facts of the initial state and of
      positive effects, the only facts a reachable state contains \\
    Thm.~4.3.8 & The formal $\tnI$ mislabels ids ($t_2$ gets two labels) & Each of the $2n+1$
      actions occurs once \\
    Thm.~4.3.8, 4.3.12 & Variables numbered $x_1, \dots, x_n$ & Variables are arbitrary numbers in
      the code of a formula; the encoding renumbers the distinct ones (otherwise it is not
      polynomial) \\
    Thm.~4.3.12 & Breaks for formulas without variables (no initial task $x_1$)
      & A dummy variable is added \\
    Thm.~4.1.1 & Assumes Chomsky normal form & Not needed; the encoding handles arbitrary rules
      including $\varepsilon$-rules \\
    \bottomrule
  \end{tabularx}
\end{table}

\section{Proofs That Take a Different Route}\label{sec:routes}

\paragraph{Lemma 4.2.11 ($\cls{ACYC} \cap \cls{REG} \in \cls{NP}$).} The article argues that
a regular problem has height~1 by Theorem~3.0.1, so Corollary~4.2.8 applies. But
Theorem~3.0.1 is about $\le_r$-stratifications, while Corollary~4.2.8 uses the height of a
$\le_0$-stratification. A $\le_0$-stratification of an acyclic problem is a total order of its
reachable compound tasks, so its height is their number. For example, the chain $A_1 \to a\,A_2
\to \dots \to a\,A_n$ is acyclic and regular, but has height $n$. The claim itself is true. The
formal proof (\lean{size_le_of_reg}) shows that decomposition trees of regular acyclic problems
are caterpillars of polynomial size, which gives polynomial certificates directly.

\paragraph{Theorem 4.2.5 ($\cls{ACYC} \in \cls{NEXPTIME}$).} The article's guess-and-check
argument is kept, but it bounds the certificate size by Lemma~4.2.4, whose bound is false. The
formal proof bounds the certificate directly: a decomposition tree rooted at a task of rank $r$
has at most $(B+1)^{r+1}$ nodes, where $B$ bounds method sizes (\lean{size_le_of_acycCert}).

\paragraph{Corollary 4.2.23 ($\cls{UNRES} \cap \cls{TO-M}$ under task insertion is in
\cls{PSPACE}).} The article restricts to ``the acyclic variant'' and uses the progression bounds
of Lemmas~4.2.15--4.2.21. Those bounds rely on a stratification, which unrestricted problems do
not have, and the restriction is not made precise. We prove a new bound
(\lean{tihtn_bound_TOM}). Acyclic progression on a problem is ordinary progression on the
problem unfolded by ancestor sets. That unfolding is acyclic, with rank $2(N - |S|) + [c \in S]$
for a task $c$ with ancestor set $S$ among $N$ compound tasks. The resulting progression bound is
$k \cdot (1 + (m-1)(2N+2))$, which is polynomial.

\paragraph{Theorem 4.3.6 (\cls{TO} is \cls{EXPTIME}-hard).} The article reduces from the game
$G_6$ and cites its \cls{EXPTIME}-hardness \parencite{stockmeyer:chandra:siamjc-79}, which is
not formalised. The formal proof uses the same technique on a different source problem. Methods
whose actions are followed by inverse actions restore the state, so a sequence of tasks acts as
a universal branching, while the choice of a method acts as an existential one. This is
applied to the tableau check of an alternating characterisation of \cls{EXPTIME}
\parencite{chandra:et:al:jacm-81}. A deterministic machine decides an exponentially padded
language of \cls{P} in exponential time, and a task $\mathit{chk}(a)$ checks the tableau entry at
time $t$ and cell $i$ using binary counters of polynomially many bits
(\lean{HTN.Article.TabData.solvable_iff}).

\paragraph{Theorem 4.3.7, acyclic part (\cls{NEXPTIME}-hard).} The printed construction is
unsound: the head counter runs one unit ahead of the tape, and the head positions count up to
$K + 2$ instead of $K$. The formal proof keeps the article's approach with a repaired
construction (\lean{TMEnc.problemFix}). The simulated machine is first made to return its head
to the start before accepting (a ``homing'' machine).

\paragraph{Theorem 4.3.7, tail-recursive part (\cls{EXPSPACE}-hard).} The article extends the
acyclic encoding by recursion but does not spell out the construction precisely enough to be
formalised. We use the tail-recursive encoding that was already proved correct in the
development. It has two alternating loops, one round per machine step; tape cells are stored in
the tails of words, and addresses are checked by binary counters.

\paragraph{Theorem 4.2.10 ($\cls{noHTN} \cap \cls{TO} \in \cls{L}$).} The article's idea is
kept, but two steps are left implicit. First, the proof iterates over the tasks of $\tnI$ in
their order, but the input lists the order only as a relation. The formal proof finds the task
of rank $k$, that is, the one with $k$ predecessors, by counting. Second, the article does not
discuss how to recognise, in logarithmic space, that the input encodes a well-formed problem of
the class. The formal proof parses codes by positions: a value nested inside $d$
self-delimiting strings sits in the input at stride $2^d$. Finally, execution is checked per
occurrence of a fact in a precondition or in the goal, rather than per state variable. This is
the same idea, with one bit of state.

\paragraph{Theorem 4.1.1 (\cls{UNRES} undecidable).} The reduction is the article's. The cited
undecidability of the intersection of context-free languages is proved from two-stack machines,
via the classical construction with triples.

\section{Status}\label{sec:status}

Table~\ref{tab:status} lists every numbered result with its Lean name in
\lean{HTN.Article}. ``Article'' means the proof follows the article; ``corrected'' and ``other
route'' refer to Sections~\ref{sec:statements} and~\ref{sec:routes}.

\begin{table}
  \caption{Status of all results. Lean names are relative to \lean{HTN.Article}, some in sub-namespaces such as \lean{Problem}.}
  \label{tab:status}
  \small
  \begin{tabularx}{\linewidth}{@{}lLl@{}}
    \toprule
    Result & Lean & Route \\
    \midrule
    Prop.~2.1.1 & \lean{prop_2_1_1} & article; up to ids \\
    Thm.~3.0.1 & \lean{Problem.thm_3_0_1} & article \\
    Thm.~4.1.1, Cors.~4.1.2--4.1.3 & \lean{thm_4_1_1}, \lean{cor_4_1_2}, \lean{cor_4_1_3} & article \\
    Thm.~4.1.4 & \lean{thm_4_1_4} & article \\
    Thm.~4.2.1 & \lean{thm_4_2_1} & article \\
    Lem.~4.2.2, Cor.~4.2.3 & \lean{Problem.lem_4_2_2}, \lean{cor_4_2_3} & corrected \\
    Lem.~4.2.4 & \lean{Problem.lem_4_2_4} & corrected \\
    Thm.~4.2.5 & \lean{thm_4_2_5} & other route \\
    Thm.~4.2.6, Cor.~4.2.7 & \lean{thm_4_2_6}, \lean{cor_4_2_7} & article \\
    Cor.~4.2.8, 4.2.9 & \lean{cor_4_2_8}, \lean{cor_4_2_9} & article \\
    Thm.~4.2.10 & \lean{thm_4_2_10} & article, gaps filled \\
    Lem.~4.2.11 & \lean{lem_4_2_11}, \lean{lem_4_2_11_TO} & other route \\
    Thm.~4.2.12, Cor.~4.2.13 & \lean{thm_4_2_12_htn}/\lean{_tihtn}, \lean{cor_4_2_13_htn}/\lean{_tihtn} & article \\
    Prop.~4.2.14 & \lean{prop_4_2_14} & article \\
    Lems.~4.2.15, 4.2.19 & \lean{lem_4_2_15}, \lean{lem_4_2_19} & corrected (4.2.15) \\
    Lem.~4.2.21 & \lean{lem_4_2_21_article} & corrected \\
    Cors.~4.2.16--4.2.22 & \lean{cor_4_2_16} \dots\ \lean{cor_4_2_22} & article \\
    Cor.~4.2.23 & \lean{cor_4_2_23_tr}, \lean{_acyc}, \lean{_unres} & other route (\cls{UNRES}) \\
    Cors.~4.2.24, 4.2.25 & \lean{cor_4_2_24}, \lean{cor_4_2_25} (variants) & article \\
    Thm.~4.3.1, Cors.~4.3.2--4.3.3 & \lean{thm_4_3_1}, \lean{cor_4_3_2}, \lean{cor_4_3_3} & article; 4.3.3 for $h \ge 1$ \\
    Thms.~4.3.4, 4.3.5 & \lean{thm_4_3_4}, \lean{thm_4_3_5} & article \\
    Thm.~4.3.6 & \lean{thm_4_3_6} & other source problem \\
    Thm.~4.3.7 & \lean{thm_4_3_7_acyc}, \lean{thm_4_3_7_tr} & repaired / other route \\
    Thm.~4.3.8, Cors.~4.3.9--4.3.11 & \lean{thm_4_3_8}, \lean{cor_4_3_9} \dots\ \lean{cor_4_3_11} & article \\
    Thm.~4.3.12 & \lean{thm_4_3_12_TO}, \lean{thm_4_3_12_TOM} & article \\
    \bottomrule
  \end{tabularx}
\end{table}

\printbibliography

\end{document}
